% =====================================================================
\part{Das Verfahren}
% =====================================================================

\chapter{Anforderungen und Konformitätsmatrix}
\label{chap:anforderungen}

\section{Was „betriebssystemunabhängig“ bedeutet}

Der Entwurf verlangt, dass die Ladeprozedur \emph{betriebssystemunabhängig durch Abstraktion und Standardisierung} wird. Als Aussage über Programme heißt das: Zwei Plattformen dürfen sich nur dort unterscheiden, wo es ihr Adapter-Orakel tut.

\begin{definition}[Beobachtung und Adapter-Parametrizität]
	\label{def:parametrizitaet}
	Die \emph{Beobachtung} einer Operationsfolge ist die Folge der Rückgaben (Handle-Identitäten bis auf Umbenennung, Statuscodes, Felder \texttt{failed\_name}, \texttt{searched}, \texttt{native\_error}) zusammen mit dem Ablauf der Hook-Aufrufe. Ein Lader heißt \emph{adapter-parametrisch}, wenn die Beobachtung bei gleicher Operationsfolge nur von $(E,R,A,h,\mathcal{F})$ und den Antworten des Orakels $\mathcal{O}$ abhängt, nicht aber vom Betriebssystem, auf dem er läuft.
\end{definition}

Satz~\ref{satz:parametrisch} zeigt, in welchem Sinn und mit welchen Ausnahmen die Referenzimplementierung diese Eigenschaft hat.

\section{Konformitätsmatrix}

Tabelle~\ref{tab:konformitaet} stellt jede Forderung des Entwurfs (Abschnittsnummern nach \texttt{proposal.md}) dem Quelltext gegenüber. „Erfüllt“ heißt: aus dem Quelltext gelesen und, wo möglich, in einem Satz bewiesen. Die Befunde sind in Kapitel~\ref{chap:befunde} ausgeführt.

{\small
\begin{longtable}{@{}L{4.1cm}L{4.5cm}L{2.2cm}L{3.6cm}@{}}
\caption{Forderungen des Entwurfs und ihre Umsetzung in \texttt{sl}.}\label{tab:konformitaet}\\
\toprule
Forderung (Abschnitt) & Umsetzung & Status & Bezug \\
\midrule
\endfirsthead
\toprule
Forderung (Abschnitt) & Umsetzung & Status & Bezug \\
\midrule
\endhead
\bottomrule
\endfoot
Namensauflösung, \texttt{.sl} anhängen, keine Plattformaffixe (5.3/1) & \texttt{search.resolve\_name} & teilweise & Lemma~\ref{lem:name}; \texttt{lib}-Präfix als Bezeichner zulässig (B14) \\[0.3em]
Ortung, alle geprüften Kandidaten aufbewahren (5.3/2) & \texttt{search.locate}, Liste \texttt{tried} & teilweise & nur im Fehlerfall zurückgegeben (B14) \\[0.3em]
Containervalidierung vor dem Abbilden (5.3/3) & \texttt{loader.\_load} & erfüllt & Satz~\ref{satz:erfolg} \\[0.3em]
Abhängigkeiten rekursiv, Zyklen erkennen (5.3/4) & \texttt{loader.\_open\_dependency}, \texttt{\_loading} & erfüllt & Satz~\ref{satz:erfolg} \\[0.3em]
Abbilden über Plattformadapter (5.3/5) & \texttt{adapters.map\_payload} über temporäre Datei & erfüllt (Stufe~1) & Kapitel~\ref{chap:adapter} \\[0.3em]
Init: Abhängigkeiten zuerst, einmal; \texttt{DllMain}/Konstruktoren umhüllt (5.3/6) & \texttt{loader.\_run\_init}; Konstruktoren laufen im Systemlader & teilweise & Satz~\ref{satz:reihenfolge}; B14 \\[0.3em]
Exporttabelle, nur deklarierte Symbole (5.3/7) & \texttt{Loader.symbol} & erfüllt (API-Ebene) & Satz~\ref{satz:export} \\[0.3em]
Referenzgezählte Handles, gleiche Instanz (5.3/8) & \texttt{Handle.refcount}, \texttt{\_loaded} & erfüllt & Satz~\ref{satz:refcount}; B13 \\[0.3em]
Freigabe in umgekehrter Abhängigkeitsordnung (5.3/9) & \texttt{loader.\_release} & erfüllt & Satz~\ref{satz:reihenfolge} \\[0.3em]
Suchordnung: explizit, \texttt{SL\_PATH}, Aufrufer, Anwendung, System (5.4) & \texttt{search.candidates} & erfüllt; App-Pfad nur falls gesetzt & Satz~\ref{satz:struktur} \\[0.3em]
Ein einheitliches Trennzeichen mit Maskierung (5.4) & \texttt{search.split\_sl\_path} & abweichend & Satz~\ref{satz:split}; B1 \\[0.3em]
Arbeitsverzeichnis, \texttt{PATH}, \texttt{LD\_LIBRARY\_PATH} nie (5.4) & \texttt{search.candidates} & erfüllt unter Bedingung & Satz~\ref{satz:cwd}; B2 \\[0.3em]
API \texttt{sl\_open}, \texttt{sl\_symbol}, \texttt{sl\_info}, \texttt{sl\_close} (5.5) & \texttt{capi}, \texttt{binding} & erfüllt (Python); C-ABI des Laders fehlt & Kapitel~\ref{chap:adapter} \\[0.3em]
Fehlermodell: Codes, \texttt{failed\_name}, \texttt{searched}, \texttt{native\_error} (5.6) & \texttt{errors.SLError}, \texttt{capi.ErrorInfo} & erfüllt; mehrzeiliger Text fehlt & B14 \\[0.3em]
\texttt{open}/\texttt{close} threadsicher (5.7) & \texttt{RLock} & erfüllt & Satz~\ref{satz:nutzung} \\[0.3em]
Init nie unter einer Sperre, die Nutzercode betreten kann (5.7) & \texttt{RLock} wird gehalten & nicht erfüllt & Satz~\ref{satz:deadlock}; B4 \\[0.3em]
Zyklus bei \texttt{open} im Init als \texttt{CYCLE} (5.7) & \texttt{\_loading} & erfüllt & Satz~\ref{satz:erfolg} \\[0.3em]
Export standardmäßig verborgen, ein Makro (5.8) & \texttt{SL\_EXPORT}, \texttt{-fvisibility=hidden} & erfüllt & Kapitel~\ref{chap:codegen} \\[0.3em]
Konformitätssuite (5.9) & Unit-Tests mit \texttt{FakeAdapter} & teilweise & Kapitel~\ref{chap:berichtet} \\[0.3em]
Metadatenfelder (6.2) & \texttt{metadata.Metadata} & erfüllt; \texttt{runtime} ohne Wirkung & B11 \\[0.3em]
Echter Container, Fat-Binary (6.3) & \texttt{container} & erfüllt & Kapitel~\ref{chap:container} \\[0.3em]
Legacy-Modus für \texttt{.dll}/\texttt{.so}/\texttt{.dylib} (8.3) & -- & fehlt & B14 \\[0.3em]
Hash und Signatur (10) & SHA-256 je Payload, keine Signatur & teilweise & Satz~\ref{satz:integritaet}; B7 \\[0.3em]
Gehärteter Modus, fail closed (10) & \texttt{hardened} & teilweise & Satz~\ref{satz:gehaertet}; B3 \\[0.3em]
Bei widersprüchlichen Kandidaten Fehler (10) & erster Treffer gewinnt & abweichend & Satz~\ref{satz:auswahl} \\
\end{longtable}
}

\chapter{Das Container-Format}
\label{chap:container}

\section{Aufbau der Datei}

Ein \SL-Container besteht aus Kopf, Metadaten, Payload-Tabelle und Payload-Bereich (Abbildung~\ref{abb:container}). Die Kodierung steht in \texttt{container.py}.

\begin{lstlisting}[language=Python, caption={Layout und Kodierung (\texttt{container.py}, Z.~12--15, 54--65).}, label={lst:container}]
MAGIC = b"\x89SLIB\r\n\x1a"  # genau 8 Byte
CONTAINER_VERSION = 1
_HEADER = struct.Struct("<8sHII")       # magic, version, meta_len, entry_count
_ENTRY = struct.Struct("<16s16sQQ32s")  # platform, arch, offset, length, sha256

def pack(metadata: Metadata, payloads: dict) -> bytes:
    if set(payloads) != set(metadata.targets):
        raise _bad("payload targets differ from metadata targets")
    meta_bytes = json.dumps(metadata.to_dict(), sort_keys=True).encode("utf-8")
    table, blob = [], bytearray()
    for target in sorted(payloads):
        data = bytes(payloads[target])
        table.append(_ENTRY.pack(target.platform.encode("ascii"), target.arch.encode("ascii"),
                                 len(blob), len(data), hashlib.sha256(data).digest()))
        blob += data
    head = _HEADER.pack(MAGIC, CONTAINER_VERSION, len(meta_bytes), len(table))
    return head + meta_bytes + b"".join(table) + bytes(blob)
\end{lstlisting}

\begin{figure}[!ht]
	\centering
	\begin{tikzpicture}[font=\small, >={Stealth[length=2mm]},
		b/.style={draw=mainblue, fill=lightgray, minimum height=1.0cm, align=center, inner sep=2pt}]
		\node[b, minimum width=2.4cm] (h) at (0,0) {Kopf\\18 Byte};
		\node[b, minimum width=3.2cm, right=0pt of h] (m) {Metadaten (JSON, UTF-8)\\$m$ Byte};
		\node[b, minimum width=3.3cm, right=0pt of m] (t) {Tabelle\\$k\cdot 80$ Byte};
		\node[b, minimum width=4.6cm, right=0pt of t] (p) {Payload-Bereich\\$\beta=\sum_t |P(t)|$ Byte};
		\node[font=\scriptsize, align=left, text width=5cm, anchor=north west] at ([yshift=-0.3cm]h.south west) {Kopf: Magic (8), Version (2), $m$ (4), $k$ (4)};
		\node[font=\scriptsize, align=left, text width=6.5cm, anchor=north west] at ([yshift=-0.3cm]t.south west) {Eintrag: Plattform (16), Architektur (16), Offset (8), Länge (8), SHA-256 (32)};
	\end{tikzpicture}
	\caption{Byte-Layout eines Containers. Ein Eintrag der Tabelle hat $16+16+8+8+32=80$ Byte, der Kopf $8+2+4+4=18$ Byte (Little Endian, ohne Auffüllbytes wegen des Präfixes \texttt{<}).}
	\label{abb:container}
\end{figure}

\begin{lemma}[Größenformel]
	\label{lem:groesse}
	Für gültige $(M,P)$ mit $k=|T|$ Targets und Metadaten-Kodierung der Länge $m$ gilt
	\[
	|\mathrm{pack}(M,P)| = 18 + m + 80\,k + \sum_{t\in T}|P(t)|.
	\]
\end{lemma}

\begin{proof}
	Der Kopf hat nach der Formatzeichenkette \texttt{<8sHII} $8+2+4+4=18$ Byte, ein Tabelleneintrag nach \texttt{<16s16sQQ32s} $16+16+8+8+32=80$ Byte. \texttt{pack} legt genau $k$ Einträge an, eine Zeile je Target aus \texttt{sorted(payloads)}, und hängt die Payloads lückenlos an den Payload-Bereich.
\end{proof}

\section{Rundreise}

\begin{lemma}[Metadaten-Rundreise]
	\label{lem:metarund}
	Für jeden gültigen Metadatensatz $M$ gilt \texttt{Metadata.from\_dict(json.loads(json.dumps(M.to\_dict(), sort\_keys=True))) = M}.
\end{lemma}

\begin{proof}
	\texttt{to\_dict} erzeugt ein Objekt aus Zeichenketten, einer ganzen Zahl, Listen, Wörterbüchern und \texttt{None}. Solche Werte werden von JSON ohne Verlust abgebildet; \texttt{None} wird zu \texttt{null} und zurück, die Schlüsselreihenfolge ändert die Werte nicht. Die Felder werden in \texttt{from\_dict} einzeln zurückgewonnen: Name, Version, ABI-Version, Laufzeit, Hooks unverändert; die Targets als Tupel in derselben Reihenfolge, in der \texttt{to\_dict} sie ausgibt; \texttt{exports} als Tupel derselben Liste; \texttt{dependencies} und \texttt{signatures} als Wörterbücher mit gleicher Abbildung (die Gleichheit von Wörterbüchern ignoriert die Einfügereihenfolge). Der Konstruktor prüft dieselben Bedingungen wie beim ersten Aufbau und akzeptiert sie wieder, weil sie nur von diesen Werten abhängen. Da die Dataclass \texttt{Metadata} komponentenweise vergleicht, folgt die Gleichheit.
\end{proof}

\begin{satz}[Rundreise des Containers]
	\label{satz:rundreise}
	Sei $(M,P)$ ein Container im Sinne von Definition~\ref{def:container}. Dann liefert \texttt{unpack(pack($M,P$))} einen Container mit Metadaten $M$, und \texttt{payload\_for}$(p,a)$ gibt für jedes Target $(p,a)\in T$ genau $P(p,a)$ zurück.
\end{satz}

\begin{proof}
	Sei $d=\mathrm{pack}(M,P)$, $m$ die Länge der Metadaten-Kodierung und $k=|T|$. Der Kopf wird mit Magic, Version~1, $m$ und $k$ gepackt; \texttt{unpack} liest sie zurück, die Prüfungen auf Länge, Magic und Version bestehen. Nach Lemma~\ref{lem:groesse} ist $18+m+80k\le|d|$, die Tabelle liegt also vollständig in $d$. Die Metadaten werden nach Lemma~\ref{lem:metarund} zu $M$ zurückgelesen. Seien $t_0<\cdots<t_{k-1}$ die Targets in aufsteigender Ordnung. Eintrag $i$ trägt den Offset $o_i=\sum_{j<i}|P(t_j)|$ und die Länge $|P(t_i)|$; der Payload-Bereich \texttt{blob} ist $d[18+m+80k:]$ und hat die Länge $\beta=\sum_j|P(t_j)|$, so dass $o_i+|P(t_i)|\le\beta$ gilt. Die Namen der Targets sind ASCII und kürzer als 16 Byte (längste Namen \texttt{windows} und \texttt{aarch64} mit 7 Byte), das Auffüllen mit Nullbytes und das Abschneiden mit \texttt{rstrip} sind also umkehrbar. Die Menge der gelesenen Targets ist $T$, die Prüfung auf Gleichheit der Mengen besteht. \texttt{payload\_for} schneidet $\mathrm{blob}[o_i:o_i+|P(t_i)|]=P(t_i)$ aus und vergleicht dessen SHA-256 mit dem gespeicherten Wert, der genau dieser Hash ist.
\end{proof}

\section{Akzeptanzbedingungen}

\begin{lemma}[Prüfungen von \texttt{unpack}]
	\label{lem:akzept}
	Wenn \texttt{unpack}$(d)$ einen Container zurückgibt, dann gilt:
	\begin{enumerate}[label=(S\arabic*)]
		\item $|d|\ge 18$;
		\item die ersten 8 Byte sind das Magic;
		\item die Version im Kopf ist $1$;
		\item $18+m+80k\le|d|$ für die Werte $m,k$ des Kopfes;
		\item die Bytes $d[18:18+m]$ sind ein JSON-Dokument, aus dem \texttt{Metadata.from\_dict} einen gültigen Metadatensatz bildet;
		\item jeder der $k$ Tabelleneinträge nennt ASCII-Namen eines gültigen Targets, und es gilt $o+\ell\le\beta$ für Offset $o$ und Länge $\ell$;
		\item die Menge der Targets der Tabelle ist gleich der Targetmenge der Metadaten.
	\end{enumerate}
\end{lemma}

\begin{proof}
	Jede Bedingung entspricht einer \texttt{raise}-Anweisung in \texttt{unpack} (Zeilen 69--97): zu kurze Datei, falsches Magic, falsche Version, abgeschnittener Container, unlesbare Metadaten, Nicht-ASCII-Namen und ungültiges \texttt{Target} (die Konstruktion löst \texttt{INVALID\_CONTAINER} aus), Payload außerhalb, abweichende Targetmengen. Ein Rückgabewert entsteht nur, wenn keine dieser Anweisungen ausgeführt wurde.
\end{proof}

\section{Integrität und Kanonizität}

\begin{satz}[Integrität der Payloads]
	\label{satz:integritaet}
	Sei $f=\mathrm{pack}(M,P)$, $f'\ne f$ gleich lang, und $f'$ unterscheide sich von $f$ nur in Bytes des Payload-Bereichs, die zum Target $t$ gehören. Wenn \texttt{unpack}$(f')$ gelingt, dann wirft \texttt{payload\_for}$(t)$ auf $f'$ eine Ausnahme \texttt{INVALID\_CONTAINER}, außer es wurde eine Kollision von SHA-256 gefunden, also zwei verschiedene Byte-Folgen mit gleichem Hash.
\end{satz}

\begin{proof}
	Die Tabelle ist unverändert, der gespeicherte Hash also der von $P(t)$. Die in $f'$ ausgeschnittene Folge $P'(t)$ hat dieselbe Länge und ist von $P(t)$ verschieden. \texttt{payload\_for} vergleicht $\mathrm{SHA256}(P'(t))$ mit dem gespeicherten Wert $\mathrm{SHA256}(P(t))$ und wirft bei Ungleichheit. Gleichheit bei $P'(t)\ne P(t)$ wäre eine Kollision.
\end{proof}

\begin{bemerkung}[Was der Hash nicht schützt]
	\label{bem:hashgrenze}
	Der Hash liegt in derselben Datei wie der Payload und schützt deshalb nur gegen \emph{zufällige} Veränderung des Payloads. Ein Angreifer, der den Payload ändert, ändert auch den Hash. Metadaten, Tabelle und Kopf sind von keinem Hash erfasst: Aus \texttt{"version": "1.2.0"} wird durch Ändern eines Bytes \texttt{"1.3.0"}, und \texttt{unpack} akzeptiert die Datei. Echtheit verlangt eine Signatur (Abschnitt~\ref{sec:signatur-ausblick}).
\end{bemerkung}

\begin{satz}[Nichtkanonizität]
	\label{satz:kanon}
	\texttt{unpack} ist auf der Menge der akzeptierten Dateien nicht injektiv bis auf Beobachtung: Es gibt verschiedene Dateien $d\ne d'$, für die \texttt{unpack} Container mit gleichen Metadaten und gleichen \texttt{payload\_for}-Ergebnissen liefert. Insbesondere ist $\mathrm{pack}\circ\mathrm{unpack}\ne\mathrm{id}$.
\end{satz}

\begin{proof}
	Sei $d=\mathrm{pack}(M,P)$. Folgende Dateien werden akzeptiert und sind beobachtbar gleich:
	\begin{enumerate}[label=(\alph*)]
		\item $d'=d\,\Vert\,\texttt{00}$: der angehängte Byte liegt im Payload-Bereich hinter allen Einträgen; keine Bedingung aus Lemma~\ref{lem:akzept} verbietet überzählige Bytes.
		\item Die JSON-Kodierung der Metadaten mit anderem Leerraum oder anderer Schlüsselreihenfolge: \texttt{json.loads} liefert dasselbe Objekt; die Längenangabe $m$ im Kopf wird entsprechend angepasst.
		\item Die Tabelleneinträge in anderer Reihenfolge: \texttt{unpack} sammelt sie in einem Wörterbuch nach Target; die Reihenfolge geht nicht ein.
		\item Zwei Targets mit gleichem Payload verweisen auf denselben Bereich: (S6) fordert nur $o+\ell\le\beta$, nicht Überlappungsfreiheit, und der Hash stimmt für beide.
		\item Ein doppelt eingetragenes Target: das Wörterbuch behält den letzten Eintrag, (S7) vergleicht Mengen.
	\end{enumerate}
	In jedem Fall unterscheiden sich die Bytes von $d$, die Beobachtung aber nicht. Da $\mathrm{pack}$ nur die kanonische Form erzeugt, folgt $\mathrm{pack}(\mathrm{unpack}(d'))\ne d'$.
\end{proof}

Die Folge für eine spätere Signatur ist praktisch: Wer eine Signatur über die Datei bildet, muss entweder eine kanonische Form erzwingen oder über wohldefinierte Teile (Metadaten-Bytes, Tabelle, Payload-Hashes) signieren. Beides ist in der Referenzimplementierung nicht angelegt (Befund~\ref{bef:canon}).

\chapter{Suchordnung}
\label{chap:suche}

\section{Namensauflösung}

\begin{lemma}[Zulässige Namen]
	\label{lem:name}
	Sei $x$ eine Zeichenkette ohne \texttt{/} und \texttt{\textbackslash}. \texttt{resolve\_name}$(x)$ liefert genau dann ein Ergebnis, wenn $x\in I\cup I\cdot\texttt{.sl}$ mit $I=\{\texttt{[A-Za-z\_][A-Za-z0-9\_]*}\}$; das Ergebnis ist \texttt{(base + ".sl", False)}. Enthält $x$ ein \texttt{/} oder \texttt{\textbackslash}, so ist der Name ein expliziter Pfad: \texttt{(x', True)} mit $x'=x$, falls $x$ auf \texttt{.sl} endet, sonst $x'=x+\texttt{.sl}$. In allen anderen Fällen wird \texttt{NOT\_FOUND} geworfen.
\end{lemma}

\begin{proof}
	Ohne Trennzeichen entfernt der Code genau ein endendes \texttt{.sl} und prüft die Basis mit der regulären Ausdrucksform $I$. Eine Basis in $I$ enthält keinen Punkt, so dass \texttt{libwsat.so}, \texttt{wsat.dll} und \texttt{wsat.sl.sl} ausscheiden, \texttt{wsat} und \texttt{wsat.sl} bestehen. Die leere Zeichenkette liegt nicht in $I$. Der explizite Pfad folgt direkt aus dem ersten \texttt{if}.
\end{proof}

\begin{bemerkung}
	Plattformsuffixe (\texttt{.so}, \texttt{.dll}, \texttt{.dylib}) sind damit \emph{syntaktisch} ausgeschlossen, Plattformpräfixe aber nicht: \texttt{libwsat} ist ein Bezeichner und wird als logischer Name \texttt{libwsat.sl} gesucht. Die Forderung des Entwurfs, das Präfix sei nicht Teil des logischen Namens, ist damit eine Konvention, keine Prüfung (B14).
\end{bemerkung}

\section{Die Kandidatenliste}

\begin{lstlisting}[language=Python, caption={Kandidatenliste (\texttt{search.py}, Zeilen 97--110).}, label={lst:candidates}]
def candidates(filename, explicit, *, env, requester_dir, app_dir, adapter, hardened) -> list:
    if explicit:
        return [filename]
    out = []
    if not hardened:
        out += [_join(d, filename) for d in split_sl_path(env.get("SL_PATH", ""))
                if _is_absolute(d)]
    if requester_dir:
        out += [_join(requester_dir, filename),
                _join(requester_dir, "sl", filename)]
    if app_dir:
        out.append(_join(app_dir, filename))
    out += [_join(d, filename) for d in adapter.system_dirs(env)]
    return list(dict.fromkeys(out))
\end{lstlisting}

\begin{definition}[Kandidatenliste]
	Für Dateinamen $f$, $e=$ \texttt{explicit} und eine Konfiguration $\mathcal{K}$ sei $\mathcal{C}(f,e,\mathcal{K})$ die von Listing~\ref{lst:candidates} berechnete Liste. Mit $\mathrm{dd}(L)$ bezeichnen wir die Liste $L$ ohne spätere Wiederholungen (\texttt{dict.fromkeys}).
\end{definition}

\begin{satz}[Struktur und Umgebungsabhängigkeit]
	\label{satz:struktur}
	Für $e=\mathrm{falsch}$ gilt $\mathcal{C}(f,e,\mathcal{K})=\mathrm{dd}(L_{\mathrm{env}}\,\Vert\,L_{\mathrm{req}}\,\Vert\,L_{\mathrm{app}}\,\Vert\,L_{\mathrm{sys}})$ mit
	\begin{align*}
	L_{\mathrm{env}}&=[\,d/f : d\in\mathrm{split}(E[\texttt{SL\_PATH}]),\ d\text{ absolut}\,]\ \text{ falls }h=\mathrm{falsch},\ \text{sonst leer},\\
	L_{\mathrm{req}}&=[\,R/f,\ R/\texttt{sl}/f\,]\ \text{ falls }R\text{ gesetzt},\quad L_{\mathrm{app}}=[\,A/f\,]\ \text{ falls }A\text{ gesetzt},\\
	L_{\mathrm{sys}}&=[\,d/f : d\in\Sigma_{\mathrm{sys}}(E)\,].
	\end{align*}
	Für $e=\mathrm{wahr}$ ist die Liste $[f]$. Die Liste hängt von der Umgebung $E$ nur über $E[\texttt{SL\_PATH}]$ (bei $h=\mathrm{falsch}$) und über $\Sigma_{\mathrm{sys}}(E)$ ab. Bei den Adaptern für Linux und macOS ist $\Sigma_{\mathrm{sys}}$ konstant; beim Windows-Adapter hängt es von \texttt{ProgramData} und \texttt{ProgramFiles} ab. Von \texttt{PATH}, \texttt{LD\_LIBRARY\_PATH} und \texttt{DYLD\_LIBRARY\_PATH} hängt die Liste in keinem Fall ab.
\end{satz}

\begin{proof}
	Die Gleichung ist der Wortlaut von Listing~\ref{lst:candidates}. Auf $E$ greift der Code nur an zwei Stellen zu: \texttt{env.get("SL\_PATH", "")} im Zweig $\neg h$ und \texttt{adapter.system\_dirs(env)}. Die Methoden \texttt{LinuxAdapter.system\_dirs} und \texttt{DarwinAdapter.system\_dirs} geben konstante Listen zurück (\texttt{adapters.py} Zeilen 75 und 82), \texttt{WindowsAdapter.system\_dirs} liest \texttt{env[key]} für \texttt{ProgramData} und \texttt{ProgramFiles} (Zeile 89). Andere Schlüssel werden nicht gelesen.
\end{proof}

\begin{satz}[Freiheit vom Arbeitsverzeichnis]
	\label{satz:cwd}
	Sei $e=\mathrm{falsch}$, seien $R$ und $A$ nicht gesetzt oder absolut, und seien alle Elemente von $\Sigma_{\mathrm{sys}}(E)$ absolut. Dann ist jeder Kandidat in $\mathcal{C}(f,e,\mathcal{K})$ ein absoluter Pfad, und die Liste sowie das Ergebnis von \texttt{locate} hängen nicht vom aktuellen Arbeitsverzeichnis ab. Ohne die Voraussetzung über $R$ gilt das nicht.
\end{satz}

\begin{proof}
	Die Elemente von $L_{\mathrm{env}}$ sind nach dem Filter \texttt{\_is\_absolute} absolut. \texttt{\_join} hängt an eine absolute Basis nur Komponenten an, das Ergebnis beginnt mit der Basis und ist absolut. Für $L_{\mathrm{req}}$, $L_{\mathrm{app}}$ und $L_{\mathrm{sys}}$ gilt das nach Voraussetzung. Ein absoluter Pfad wird vom Dateisystem unabhängig vom Arbeitsverzeichnis aufgelöst. Gegenbeispiel: Für $R=\texttt{rel}$ enthält die Liste \texttt{rel/x.sl}, das relativ zum Arbeitsverzeichnis aufgelöst wird.
\end{proof}

Der Lader erzwingt die Voraussetzung nicht: \texttt{Loader.\_\_init\_\_} übernimmt \texttt{requester\_dir} und \texttt{app\_dir} unverändert (Befund~\ref{bef:cwd}). Für Abhängigkeiten ist sie erfüllt, weil das Aufruferverzeichnis aus dem Schlüssel (\texttt{realpath}) gebildet wird.

\begin{lemma}[Entdoppelung]
	\label{lem:dedup}
	$\mathrm{dd}(L)$ enthält jedes Element von $L$ genau einmal, in der Reihenfolge seines ersten Auftretens in $L$.
\end{lemma}

\begin{proof}
	\texttt{dict.fromkeys(L)} fügt die Elemente von $L$ der Reihe nach als Schlüssel ein; ein vorhandener Schlüssel wird nicht erneut eingefügt und behält seine Position. Die Sprachdefinition von Python garantiert seit Version 3.7 die Einfügereihenfolge von Wörterbüchern.
\end{proof}

\section{Auswahl des Treffers}

\begin{lstlisting}[language=Python, caption={Auswahl (\texttt{search.py}, Zeilen 130--145).}, label={lst:locate}]
def locate(name_or_path, *, env, requester_dir, app_dir, adapter, hardened=False) -> tuple:
    filename, explicit = resolve_name(name_or_path)
    tried = []
    for path in candidates(filename, explicit, env=env, requester_dir=requester_dir,
                           app_dir=app_dir, adapter=adapter, hardened=hardened):
        tried.append(path)
        if os.path.isfile(path):
            dir_to_check = os.path.dirname(os.path.abspath(path))
            if hardened and _world_writable(dir_to_check):
                raise SLError(SLStatus.ACCESS_DENIED,
                              f"'{path}' lies in a world-writable directory",
                              failed_name=name_or_path, searched=tried)
            return os.path.abspath(path), tried
    raise SLError(SLStatus.NOT_FOUND, f"library '{name_or_path}' not found",
                  failed_name=name_or_path, searched=tried)
\end{lstlisting}

\begin{satz}[Erster Treffer gewinnt]
	\label{satz:auswahl}
	Sei $\mathcal{C}=[p_1,\ldots,p_n]$ die Kandidatenliste und $i$ der kleinste Index, für den $p_i$ eine reguläre Datei ist. Dann gilt bei $h=\mathrm{falsch}$: \texttt{locate} gibt $(\mathrm{abspath}(p_i),[p_1,\ldots,p_i])$ zurück. Gibt es kein solches $i$, wird \texttt{NOT\_FOUND} mit \texttt{searched}$=[p_1,\ldots,p_n]$ geworfen. Die Zahl der Dateisystemprüfungen ist höchstens $n\le |L_{\mathrm{env}}|+3+|L_{\mathrm{sys}}|$. Spätere Kandidaten werden nach einem Treffer nicht mehr betrachtet, auch wenn der Treffer in der weiteren Prüfung (Version, ABI) scheitert.
\end{satz}

\begin{proof}
	Die Schleife durchläuft die Liste in Reihenfolge, fügt jeden Pfad vor der Prüfung an \texttt{tried} an und kehrt beim ersten Treffer zurück. Erschöpft sich die Schleife, wird mit der vollständigen Liste geworfen. Die Schranke folgt aus Satz~\ref{satz:struktur}. Die Prüfung von Version und ABI findet erst im Lader nach der Rückgabe statt; ein Fehler dort ruft \texttt{locate} nicht erneut mit späteren Kandidaten auf.
\end{proof}

Die letzte Aussage hat Folgen: Der Entwurf verlangt (Abschnitt~10) bei widersprüchlichen Kandidaten einen Fehler statt einer stillen Wahl. Die Implementierung wählt still den ersten. Das ist für die Reproduzierbarkeit gut (Ergebnis folgt der festen Rangfolge), für die Erkennung von Verwechslungen schwach; Kapitel~\ref{chap:versionen} zeigt, dass eine Suche über \emph{alle} Kandidaten mit Rückverfolgung das Problem schwer machen würde.

\section{Der gehärtete Modus}

\begin{satz}[Gehärteter Modus]
	\label{satz:gehaertet}
	Sei $h=\mathrm{wahr}$ und $e=\mathrm{falsch}$. Dann gilt:
	\begin{enumerate}[label=(\roman*)]
		\item \texttt{SL\_PATH} hat keinen Einfluss auf die Kandidatenliste.
		\item Gibt \texttt{locate} einen Pfad zurück, so hat das Verzeichnis des Treffers auf einem POSIX-System das Bit \texttt{S\_IWOTH} nicht gesetzt.
		\item Liegt der erste existierende Kandidat in einem solchen Verzeichnis, wird \texttt{ACCESS\_DENIED} geworfen, auch wenn ein späterer Kandidat in einem sicheren Verzeichnis liegt (\emph{fail closed}).
	\end{enumerate}
\end{satz}

\begin{proof}
	(i) Nach Satz~\ref{satz:struktur} ist $L_{\mathrm{env}}$ bei $h=\mathrm{wahr}$ leer. (ii) Die Rückgabe erfolgt nur, wenn die Bedingung \texttt{hardened and \_world\_writable(dir)} falsch war; unter POSIX gibt \texttt{\_world\_writable} den Wert von \texttt{st\_mode \& S\_IWOTH} zurück, bei einem Fehler von \texttt{os.stat} dagegen \texttt{False}. Im Erfolgsfall hat \texttt{stat} also entweder das Bit nicht gesetzt gefunden oder ist gescheitert. (iii) Die Prüfung erfolgt beim ersten Treffer; bei Zutreffen wird sofort geworfen, spätere Kandidaten werden nicht betrachtet.
\end{proof}

Die Grenzen dieses Satzes sind Gegenstand von Befund~\ref{bef:ww}: Geprüft wird nur das unmittelbare Verzeichnis, nur das Bit für \emph{others}, nicht Besitzer, Gruppenrechte oder übergeordnete Verzeichnisse; bei einem \texttt{stat}-Fehler wird die Datei zugelassen (fail open), und unter Windows enthält der Code eine auf einen Testnamen zugeschnittene Sonderregel (Zeilen 120--123).

\section{Zerlegung von \texttt{SL\_PATH}}

\begin{lstlisting}[language=Python, caption={\texttt{split\_sl\_path} (\texttt{search.py}, Zeilen 14--65, gekürzt).}, label={lst:split}]
def split_sl_path(raw: str) -> list[str]:
    if not raw:
        return []
    # Windows Semikolon-Trennliste unterstützen
    if ";" in raw:
        return [p for p in raw.split(";") if p]
    parts, current, i, n = [], [], 0, len(raw)
    while i < n:
        # Windows-Laufwerksbuchstabe (z. B. C:\)
        if (i == 1 and raw[0].isalpha() and raw[1] == ":"
                and (n == 2 or raw[2] in ("/", "\\"))):
            current.append(":"); i += 1; continue
        if raw[i] == "\\" and i + 1 < n:
            if raw[i + 1] == ":":  current.append(":");  i += 2; continue
            if raw[i + 1] == "\\": current.append("\\"); i += 2; continue
        if raw[i] == ":":
            token = "".join(current)
            if token: parts.append(token)
            current = []; i += 1
        else:
            current.append(raw[i]); i += 1
    if current:
        token = "".join(current)
        if token: parts.append(token)
    return parts
\end{lstlisting}

Der Entwurf verlangt ein \emph{einziges} Trennzeichen mit festgelegter Maskierung, damit \texttt{SL\_PATH} auf allen Systemen gleich gelesen wird (Abschnitt~5.4 und offene Frage in Abschnitt~14). Die Umsetzung hat zwei Betriebsarten, die der \emph{Inhalt} der Zeichenkette wählt, nicht die Plattform.

\begin{definition}[Maskierung und Verbindung]
	Für eine Zeichenkette $p$ sei $\mathrm{esc}(p)$ das Ergebnis, in dem jedes \texttt{\textbackslash} durch \texttt{\textbackslash\textbackslash} und jedes \texttt{:} durch \texttt{\textbackslash:} ersetzt wird. Für eine Liste $(p_1,\ldots,p_k)$ sei $\mathrm{join}(p_1,\ldots,p_k)=\mathrm{esc}(p_1)\,\texttt{:}\,\cdots\,\texttt{:}\,\mathrm{esc}(p_k)$.
\end{definition}

\begin{satz}[Zerlegung von \texttt{SL\_PATH}]
	\label{satz:split}
	Sei $\mathrm{split}$ die Funktion \texttt{split\_sl\_path}.
	\begin{enumerate}[label=(\alph*)]
		\item $\mathrm{split}(\mathit{raw})$ ist eine Funktion von $\mathit{raw}$ allein; sie hängt weder von der Plattform noch von der Umgebung ab, und ihre Laufzeit ist $O(|\mathit{raw}|)$.
		\item Enthält $\mathit{raw}$ ein Semikolon, so ist $\mathrm{split}(\mathit{raw})$ die Liste der nichtleeren Teile von \texttt{raw.split(";")}. Jeder Teil ist semikolonfrei und wird nicht entmaskiert; der Doppelpunkt trennt dann nicht.
		\item Seien $p_1,\ldots,p_k$ nichtleere, semikolonfreie Zeichenketten. Gilt nicht zugleich $k\ge2$, $p_1$ ist ein einzelner Buchstabe und das erste Zeichen von $\mathrm{esc}(p_2)$ ist \texttt{/} oder \texttt{\textbackslash}, so ist $\mathrm{split}(\mathrm{join}(p_1,\ldots,p_k))=(p_1,\ldots,p_k)$.
		\item Kein Eintrag, der ein Semikolon enthält, ist darstellbar: Für kein $\mathit{raw}$ enthält $\mathrm{split}(\mathit{raw})$ einen Eintrag mit Semikolon. Insbesondere ist die Abbildung $\mathrm{split}$ \emph{nicht} surjektiv auf Listen nichtleerer Zeichenketten.
		\item Dieselbe Zeichenkette kann je nach Inhalt Verschiedenes bedeuten: $\mathrm{split}(\texttt{a:b})=(\texttt{a},\texttt{b})$, aber $\mathrm{split}(\texttt{a:b;c})=(\texttt{a:b},\texttt{c})$.
	\end{enumerate}
\end{satz}

\begin{proof}
	(a) Der Code liest nur $\mathit{raw}$ und besteht aus einer Schleife, die in jedem Durchlauf mindestens ein Zeichen konsumiert. (b) Der erste \texttt{if} gibt in diesem Fall sofort zurück. (d) folgt aus (b): Teile von \texttt{raw.split(";")} enthalten kein Semikolon, und ein Semikolon in einem Eintrag lässt $\mathit{raw}$ ein Semikolon enthalten, so dass nach (b) zerlegt wird. (e) ist Anwendung von (b) und (c) mit den angegebenen Zeichenketten.

	(c) Da kein $p_i$ ein Semikolon enthält, enthält auch $\mathit{raw}=\mathrm{join}(\ldots)$ keines; der Zweig aus (b) wird nicht betreten. Jede Einheit der Kodierung ist entweder ein Zeichen, das weder \texttt{\textbackslash} noch \texttt{:} ist, ein Maskierungspaar \texttt{\textbackslash\textbackslash} oder \texttt{\textbackslash:}, oder ein Trennzeichen \texttt{:} zwischen zwei Einträgen. In $\mathrm{esc}$-Ausgaben steht jedes \texttt{\textbackslash} am Anfang eines Paares, ein einzelnes \texttt{\textbackslash} kommt nicht vor, und ein unmaskiertes \texttt{:} kommt nur als Trennzeichen vor. Wir zeigen durch Induktion über die Zahl der verarbeiteten Einheiten: Steht die Schleife an einer Einheitengrenze, so enthält \texttt{current} die entmaskierten Zeichen des laufenden Eintrags und \texttt{parts} die fertigen Einträge. Beginn: leer. Schritt: Ein gewöhnliches Zeichen wird angehängt; ein Paar \texttt{\textbackslash:} bzw. \texttt{\textbackslash\textbackslash} wird durch den zweiten bzw. dritten \texttt{if}-Zweig als \texttt{:} bzw. \texttt{\textbackslash} angehängt und überspringt zwei Zeichen; ein Trennzeichen schließt den Eintrag ab (er ist nichtleer) und beginnt einen neuen. Der Laufwerkszweig könnte nur bei $i=1$ eingreifen, wenn \texttt{raw[0]} ein Buchstabe und \texttt{raw[1]} ein \texttt{:} ist. Dann ist \texttt{raw[0]} einzelnes Zeichen vor einem unmaskierten Doppelpunkt, also $p_1$ ein einzelner Buchstabe und $k\ge2$; der Zweig greift zusätzlich nur, wenn $n=2$ (unmöglich bei $k\ge2$) oder \texttt{raw[2]}, das erste Zeichen von $\mathrm{esc}(p_2)$, \texttt{/} oder \texttt{\textbackslash} ist. Das ist durch die Voraussetzung ausgeschlossen. Am Ende hängt der Code den letzten nichtleeren Eintrag an. Also ist das Ergebnis $(p_1,\ldots,p_k)$.
\end{proof}

\begin{beispiel}[Zwei Laufwerke]\label{bsp:zweilaufwerke}
	Für $\mathit{raw}=\texttt{D:\textbackslash a:C:\textbackslash b}$ erkennt der Laufwerkszweig nur das erste Laufwerk. Die Zerlegung liefert $(\texttt{D:\textbackslash a},\texttt{C},\texttt{\textbackslash b})$ und nicht die beiden Pfade. Der zweite Laufwerksbuchstabe muss maskiert werden (\texttt{C\textbackslash:\textbackslash\textbackslash b}) oder die Zeichenkette muss ein Semikolon als Trenner verwenden.
\end{beispiel}

Der Satz zeigt, dass die README-Formulierung „Trennzeichen \texttt{:} unter POSIX oder \texttt{;} unter Windows“ die Implementierung ungenau beschreibt: Die Wahl trifft die Zeichenkette, nicht das System. Das ist gerade für die Forderung der Plattformunabhängigkeit günstig (Teil~(a)), aber nicht dokumentiert, und der Zweig mit Semikolon ist im Abdeckungsbericht als nie ausgeführt ausgewiesen (Befund~\ref{bef:sep}).

\chapter{Der Lader-Kern}
\label{chap:lader}

Der Lader-Kern setzt die neun Schritte des Entwurfs (Abschnitt~5.3) um. Listing~\ref{lst:lader} zeigt die drei Funktionen, auf denen alle folgenden Sätze beruhen.

\begin{lstlisting}[language=Python, caption={\texttt{Loader.\_open}, \texttt{\_load} und \texttt{\_release} (\texttt{loader.py}, gekürzt).}, label={lst:lader}]
def _open(self, name_or_path, spec, requester_dir) -> Handle:
    path, _tried = locate(name_or_path, env=self.env, requester_dir=requester_dir,
                          app_dir=self.app_dir, adapter=self.adapter, hardened=self.hardened)
    key = os.path.realpath(path)
    if key in self._loading:
        raise SLError(SLStatus.CYCLE, f"cyclic dependency through '{name_or_path}'",
                      failed_name=name_or_path)
    handle = self._loaded.get(key)
    if handle is not None:
        self._check_range(handle.metadata, spec)
        handle.refcount += 1
        return handle
    self._loading.append(key)
    try:
        handle = self._load(key, spec)
    finally:
        self._loading.pop()
    self._loaded[key] = handle
    return handle

def _load(self, path: str, spec) -> Handle:
    container = unpack(self._read(path))                      # Schritt 3
    meta = container.metadata
    low, high = self.abi_range
    if not low <= meta.abi_version <= high:
        raise SLError(SLStatus.ABI_MISMATCH, ...)
    payload = container.payload_for(self.platform, self.arch)
    self._check_range(meta, spec)
    deps, native = [], None
    try:
        for dep, dep_spec in meta.dependencies.items():       # Schritt 4
            deps.append(self._open_dependency(meta, dep, dep_spec, path))
        native = self.adapter.map_payload(payload, meta.name)  # Schritt 5
        handle = Handle(meta.name, path, meta, native, self, deps)
        if meta.init:                                         # Schritt 6
            self._run_init(handle)
        return handle
    except SLError:
        if native is not None:
            self.adapter.unmap(native)
        for dep_handle in deps:
            self._release(dep_handle)
        raise

def _release(self, handle: Handle) -> None:
    handle.refcount -= 1
    if handle.refcount > 0:
        return
    handle.closed = True
    del self._loaded[handle.path]
    if handle.metadata.fini:
        try:
            self.adapter.lookup(handle.native, handle.metadata.fini)()
        except SLError:
            pass
    self.adapter.unmap(handle.native)
    for dep_handle in handle.deps:
        self._release(dep_handle)
\end{lstlisting}

\section{Kosten und Charakterisierung des Erfolgs}

Im Folgenden gelte Annahme~\ref{ann:determinismus}, der Anfangszustand sei $s_0=(\emptyset,[\,])$, und alle Hooks und Orakelaufrufe seien deterministisch.

\begin{definition}[Lokale Bedingung]
	\label{def:lokal}
	Eine Datei $f\in V_r$ erfüllt die \emph{lokale Bedingung}, wenn (1) $f$ ein gültiger Container ist (Lemma~\ref{lem:akzept}), (2) die ABI-Version von $M_f$ im Bereich des Laders liegt, (3) ein Payload für das Target des Laders vorhanden ist und seinen Hash besteht, (4) das Orakel den Payload abbildet, (5) ein deklarierter Init-Hook gefunden wird und $0$ liefert.
\end{definition}

\begin{satz}[Aufrufzahlen]
	\label{satz:kosten}
	Gelingt \texttt{open}$(r)$ von $s_0$ aus, so wird \texttt{\_load} genau $|V_r|$-mal und \texttt{\_open} genau $1+|E_r|$-mal aufgerufen; die Zahl der Dateisystemprüfungen ist höchstens $(1+|E_r|)\cdot n_{\max}$ mit der Kandidatenschranke $n_{\max}$ aus Satz~\ref{satz:auswahl}. Die Rekursionstiefe ist höchstens $|V_r|$.
\end{satz}

\begin{proof}
	Im Erfolgsfall durchläuft der Code jede Abhängigkeitsdeklaration jedes geladenen Knotens genau einmal und ruft für sie \texttt{\_open} auf; dazu kommt der äußere Aufruf: $1+|E_r|$ Aufrufe. Für jeden Knoten $f\in V_r$ wird \texttt{\_load} beim ersten Erreichen aufgerufen; jedes spätere Erreichen findet $f$ in $\Lambda$, weil der Knoten nach seinem Abschluss eingetragen ist und ein Erreichen \emph{während} des Ladens in \texttt{CYCLE} enden würde, also nicht im Erfolgsfall liegt. Jeder Aufruf von \texttt{\_open} führt höchstens $n_{\max}$ Prüfungen aus. Auf dem Keller stehen paarweise verschiedene Schlüssel (sonst \texttt{CYCLE}), also höchstens $|V_r|$.
\end{proof}

\begin{lemma}[Abschlussreihenfolge]
	\label{lem:abschluss}
	Gelingt \texttt{open}$(r)$, und ist $f\to g$ eine Kante von $G_r$, so ist \texttt{\_load}$(g)$ (oder, falls $g$ schon geladen war, dessen früherer Abschluss) beendet, bevor \texttt{\_load}$(f)$ beendet ist.
\end{lemma}

\begin{proof}
	\texttt{\_load}$(f)$ ruft in der Schleife \texttt{\_open} für $g$ auf, bevor es zurückkehrt. Dieser Aufruf kehrt im Erfolgsfall zurück; entweder lädt er $g$ vollständig (\texttt{\_load}$(g)$ endet vor der Rückkehr) oder findet $g$ in $\Lambda$, wo es nach Abschluss eingetragen wurde.
\end{proof}

\begin{satz}[Charakterisierung des Erfolgs; Zyklen]
	\label{satz:erfolg}
	Seien $s_0$ und Annahme~\ref{ann:determinismus} gegeben. \texttt{open}$(r)$ gelingt genau dann, wenn
	\begin{enumerate}[label=(\roman*)]
		\item der Graph $G_r$ azyklisch ist,
		\item jede Datei $f\in V_r$ die lokale Bedingung erfüllt,
		\item jeder Abhängigkeitsname jeder Datei $f\in V_r$ zu einer Datei $g=\mathrm{res}(n,\mathrm{dirname}(f))\ne\bot$ aufgelöst wird (im gehärteten Modus: nicht in einem schreibbaren Verzeichnis) und die Version von $g$ den Bereich $\mathsf{s}$ der Deklaration erfüllt.
	\end{enumerate}
	Außerdem gilt: Wird \texttt{CYCLE} geworfen, so enthält $G_r$ einen Kreis.
\end{satz}

\begin{proof}
	\emph{Notwendigkeit.} Gelingt \texttt{open}$(r)$, so wurde jeder Knoten von $V_r$ geladen und jede Deklaration abgearbeitet. Jede der Bedingungen (2) und (3) entspricht einer \texttt{raise}-Anweisung, die nicht ausgeführt wurde. Zu (i): Angenommen $G_r$ hätte einen Kreis $f_1\to\cdots\to f_j\to f_1$. Nach Lemma~\ref{lem:abschluss} endet \texttt{\_load}$(f_{i+1})$ vor \texttt{\_load}$(f_i)$ und \texttt{\_load}$(f_1)$ vor \texttt{\_load}$(f_j)$, also läge der Abschluss von $f_1$ echt vor sich selbst; ein Widerspruch. (Ein Kreis der Länge~1 endet unmittelbar in \texttt{CYCLE}, weil $f_1$ auf dem Keller steht.)

	\emph{Hinlänglichkeit.} Da $G_r$ azyklisch und endlich ist, hat jeder Knoten eine endliche Höhe $\eta(f)$, die Länge des längsten von $f$ ausgehenden Weges. Wir zeigen durch Induktion über $\eta(f)$: \texttt{\_open}$(f)$ gelingt in jedem Zustand, in dem kein Schlüssel aus $V_f$ auf dem Keller steht. Dies ist für $f=r$ im Zustand $s_0$ erfüllt, und es bleibt beim rekursiven Aufruf erhalten: Stünde ein Knoten $u\in V_g$ auf dem Keller, so wäre $u$ ein Vorfahre von $g$ auf dem Aufrufweg, hätte also einen Weg zu $g$ und läge in $V_g$; mit dem Weg von $g$ zu $u$ entstünde ein Kreis. Für $\eta(f)=0$ hat $f$ keine Abhängigkeiten; die lokale Bedingung sichert alle Schritte, $f$ steht nicht auf dem Keller, \texttt{\_open} gelingt (oder findet $f$ in $\Lambda$ und besteht die Bereichsprüfung nach (iii)). Für $\eta(f)>0$ haben alle Abhängigkeiten $g$ kleinere Höhe und erfüllen die Voraussetzung; nach Induktion und (iii) gelingen ihre Aufrufe, danach gelingen Abbilden und Init nach (2).

	\emph{Zusatz.} \texttt{CYCLE} wird nur geworfen, wenn der Schlüssel schon auf dem Keller steht. Der Keller ist zu jedem Zeitpunkt ein Weg von $r$ aus in $G_r$: Jeder Eintrag kommt aus einem \texttt{\_open}, das aus der Schleife des darunterliegenden \texttt{\_load} aufgerufen wurde, also über eine Kante. Der Weg vom früheren Eintrag zum jetzigen plus die Kante zurück zum früheren Eintrag bildet einen Kreis.
\end{proof}

Bemerkung zur Reichweite: Welcher Fehler gemeldet wird, wenn mehrere Bedingungen verletzt sind, hängt von der Reihenfolge der Schleife über \texttt{meta.dependencies.items()} ab und ist von Satz~\ref{satz:erfolg} nicht erfasst. Entscheidend ist, dass \emph{Erfolg} reihenfolgeunabhängig charakterisiert ist.

\section{Operationsvertrag, Referenzzählung und Reihenfolgen}

Für die folgenden Sätze gelte zusätzlich:

\begin{annahme}[Ausnahmeklasse]
	\label{ann:exc}
	Alle Aufrufe des Orakels und alle Hooks kehren normal zurück oder werfen eine \texttt{SLError}. Es werden keine anderen Ausnahmen geworfen.
\end{annahme}

Die Annahme ist nicht selbstverständlich; unter Windows löst ein Speicherzugriffsfehler in nativem Code in \texttt{ctypes} eine \texttt{OSError} aus. Satz~\ref{satz:ausnahme} zeigt, was geschieht, wenn sie verletzt ist.

\begin{lemma}[Operationsvertrag]
	\label{lem:vertrag}
	Sei $s=(\Lambda,\sigma)$ ein Zustand, in dem kein Schlüssel von $\sigma$ in $\mathrm{dom}\,\Lambda$ liegt und $\mathrm{dom}\,\Lambda$ unter Abhängigkeiten abgeschlossen ist. Ein Aufruf $g=\texttt{\_open}(n,\mathsf{s},d)$ verhält sich so:
	\begin{enumerate}[label=(\alph*)]
		\item \emph{Fehler.} Wirft er eine \texttt{SLError}, so ist der Zustand danach gleich $s$ (gleiches $\Lambda$, gleicher Keller, gleiche Zähler aller bereits vorhandenen Handles), und jedes während des Aufrufs erfolgte \texttt{map} ist durch ein \texttt{unmap} beantwortet, jedes erfolgreiche Init durch ein Fini.
		\item \emph{Erfolg.} Gelingt er, so ist $N$ die Menge der dabei neu geladenen Schlüssel, $\Lambda'=\Lambda\cup N$, und für jedes Handle $h$ gilt
		\[
		\mathrm{rc}'(h)=\mathrm{rc}(h)+\bigl|\{(f,i): f\in N,\ \mathrm{deps}(f)[i]=h\}\bigr|+[h=g],
		\]
		wobei $\mathrm{rc}(h)=0$ für $h\in N$ vor dem Aufruf gesetzt ist. Der Keller ist unverändert, $\mathrm{dom}\,\Lambda'$ ist wieder unter Abhängigkeiten abgeschlossen.
	\end{enumerate}
\end{lemma}

\begin{proof}
	Induktion über die Zahl der rekursiven \texttt{\_open}-Aufrufe, die der Aufruf auslöst. \emph{Fall \texttt{locate} wirft} oder \emph{Schlüssel auf dem Keller:} Es wurde nichts verändert (a). \emph{Fall Treffer in $\Lambda$:} Besteht die Bereichsprüfung nicht, wird vor der Erhöhung geworfen und nichts verändert (a); sonst wird $\mathrm{rc}(g)$ um eins erhöht, $N=\emptyset$ (b). \emph{Fall neuer Schlüssel:} Er wird auf den Keller gelegt und am Ende durch \texttt{finally} wieder entfernt, in beiden Ausgängen. Die Schritte vor dem \texttt{try} in \texttt{\_load} (Entpacken, ABI, Payload, Bereich) verändern keinen Zustand; werfen sie, ist (a) erfüllt. Im \texttt{try}: Die Aufrufe für die Abhängigkeiten gelingen nach Induktionsvoraussetzung (b) oder lassen den Zustand nach (a) unverändert. Wirft das $j$-te, so sind die Handles $\mathrm{deps}[1..j-1]$ in \texttt{deps} eingetragen; der \texttt{except}-Zweig ruft \texttt{\_release} für jedes von ihnen. Jedes davon nimmt die durch seinen Aufruf erhaltene Referenz zurück: War das Handle schon in $\Lambda$, sinkt $\mathrm{rc}$ auf den früheren Wert; war es neu, sinkt $\mathrm{rc}$ auf~$0$, es wird aus $\Lambda$ entfernt, sein Fini-Hook läuft, es wird freigegeben, und rekursiv werden seine Abhängigkeiten zurückgegeben; nach Induktion über die Höhe stellt das genau den früheren Zustand her. Wirft das Abbilden, ist \texttt{native} nicht gesetzt; wirft der Init-Hook, wird die native Ressource freigegeben (\texttt{unmap}). In beiden Fällen sind danach alle Referenzen der Abhängigkeiten zurückgegeben, das neue Handle wurde nie in $\Lambda$ eingetragen. Das ist (a); das Fini fehlt für den eigenen gescheiterten Init, weil dieser nicht gelungen ist. Gelingen alle Schritte, wird der neue Schlüssel nach dem Pop in $\Lambda$ eingetragen; jedes Handle $h$ hat für jeden Eintrag in \texttt{deps} eines neu geladenen $f$ eine Referenz erhalten (aus dem Aufruf von \texttt{\_open}, der $h$ lieferte), und das neue Handle selbst beginnt mit $\mathrm{rc}=1=[h=g]$. Abgeschlossenheit folgt, weil jede Abhängigkeit des neuen Knotens ihrerseits in $\Lambda'$ liegt.
\end{proof}

\begin{definition}[Externe Referenzen]
	Für ein Handle $h$ sei $\mathrm{ext}(h)$ die Zahl der erfolgreichen äußeren \texttt{open}-Aufrufe, die $h$ geliefert haben und denen noch kein äußeres \texttt{close}$(h)$ gegenübersteht.
\end{definition}

\begin{satz}[Invariante der Referenzzählung]
	\label{satz:refcount}
	Zwischen zwei äußeren Operationen gilt für jedes $h\in\mathrm{dom}\,\Lambda$
	\[
	\mathrm{rc}(h)=\mathrm{ext}(h)+\bigl|\{(f,i): f\in\mathrm{dom}\,\Lambda,\ \mathrm{deps}(f)[i]=h\}\bigr|.
	\]
	Insbesondere ist $\mathrm{rc}(h)\ge1$ für jedes geladene Handle. Werden alle äußeren Referenzen zurückgegeben ($\mathrm{ext}\equiv0$), so ist danach $\Lambda=\emptyset$, und jedes \texttt{map} wurde durch genau ein \texttt{unmap} beantwortet.
\end{satz}

\begin{proof}
	Induktion über die Folge der äußeren Operationen; zu Beginn ist $\Lambda$ leer. Ein äußeres \texttt{open} gelingt oder scheitert. Scheitert es, bleibt nach Lemma~\ref{lem:vertrag}(a) alles unverändert. Gelingt es, so liefert Lemma~\ref{lem:vertrag}(b): $\mathrm{rc}$ wächst für jedes Handle um die Zahl der neuen Abhängigkeitseinträge, die darauf zeigen (dieselbe Zahl, um die die zweite Summe wächst, weil genau die neu geladenen $f\in N$ hinzukommen), und das gelieferte Handle erhält einen Zähler mehr, dem $\mathrm{ext}$ um eins mehr entspricht. Ein äußeres \texttt{close}$(h)$ ruft \texttt{\_release}$(h)$. Bleibt $\mathrm{rc}>0$, sinkt $\mathrm{rc}(h)$ und $\mathrm{ext}(h)$ um eins. Wird $\mathrm{rc}=0$, so ist $\mathrm{ext}(h)=0$ und es gibt keinen Eintrag in einem anderen geladenen Handle auf $h$; $h$ wird entfernt, und für jeden seiner Abhängigkeitseinträge sinkt der Zähler des Ziels um eins, was der Verkleinerung der zweiten Summe um genau diesen Eintrag entspricht; die Invariante bleibt (rekursiv) erhalten. Sind alle $\mathrm{ext}$ null, so folgt $\mathrm{rc}(h)=$ Zahl der Abhängigkeitseinträge auf $h$. Ein geladenes Handle mit $\mathrm{rc}\ge1$ besitzt dann einen Vorgänger in $\Lambda$. Wegen der Abgeschlossenheit und der Endlichkeit führt die Vorgängerrelation zu einem Handle ohne Vorgänger oder zu einem Kreis; Kreise gibt es nach Satz~\ref{satz:erfolg} nicht. Ein Handle ohne Vorgänger hätte $\mathrm{rc}=0$, im Widerspruch zu $\mathrm{rc}\ge1$. Also ist $\Lambda=\emptyset$. Da jedes erfolgreiche \texttt{map} an ein Handle in $\Lambda$ gebunden ist, das bei seiner Entfernung genau einmal \texttt{unmap} auslöst, und gescheiterte Ladevorgänge nach Lemma~\ref{lem:vertrag}(a) ihr \texttt{map} zurücknehmen, gilt die Bilanz.
\end{proof}

\begin{satz}[Reihenfolge von Init und Fini]
	\label{satz:reihenfolge}
	Sei $f\to g$ eine Kante von $G_r$ in einem erfolgreichen \texttt{open}$(r)$. Dann gilt:
	\begin{enumerate}[label=(\roman*)]
		\item $\mathrm{init}(g)\prec\mathrm{init}(f)$; $\mathrm{map}(g)\prec\mathrm{map}(f)$.
		\item Für jede Datei läuft Init höchstens einmal pro \emph{Ladeinstanz}, also zwischen zwei \texttt{map}/\texttt{unmap}-Paaren; mehrfaches \texttt{open} derselben Datei löst Init nicht erneut aus.
		\item Werden später alle Referenzen zurückgegeben, so gilt $\mathrm{fini}(f)\prec\mathrm{fini}(g)$ und $\mathrm{unmap}(f)\prec\mathrm{unmap}(g)$.
		\item Jedes Fini folgt auf ein erfolgreiches Init derselben Ladeinstanz; die Zahl der Fini-Aufrufe ist gleich der Zahl erfolgreicher Init-Aufrufe.
	\end{enumerate}
\end{satz}

\begin{proof}
	(i) \texttt{map}$(f)$ und Init$(f)$ stehen in \texttt{\_load}$(f)$ \emph{nach} der Schleife über die Abhängigkeiten. Nach Lemma~\ref{lem:abschluss} ist \texttt{\_load}$(g)$ beendet, wenn die Schleife $g$ abgearbeitet hat, oder $g$ war früher fertig; map und Init von $g$ liegen in \texttt{\_load}$(g)$ vor dessen Ende. (ii) Init steht nur in \texttt{\_load}, das nur für Schlüssel außerhalb von $\Lambda$ aufgerufen wird; ein Schlüssel kommt erst nach \texttt{unmap} wieder aus $\Lambda$. (iii) Nach Satz~\ref{satz:refcount} enthält $\mathrm{rc}(g)$ eine Referenz für den Abhängigkeitseintrag in $f$. In \texttt{\_release}$(f)$ stehen Fini und \texttt{unmap} von $f$ \emph{vor} der Schleife, die $g$ zurückgibt; $g$ erreicht $\mathrm{rc}=0$ frühestens nach dieser Rückgabe, und Fini und \texttt{unmap} von $g$ folgen danach. (iv) Fini steht in \texttt{\_release} nur für Handles, die vollständig geladen und damit nach \texttt{\_run\_init} (soweit deklariert) in $\Lambda$ eingetragen wurden; ein gescheiterter Init erreicht $\Lambda$ nie (Lemma~\ref{lem:vertrag}(a)).
\end{proof}

\begin{satz}[Ausnahmesicherheit und ihre Grenze]
	\label{satz:ausnahme}
	(a) Unter Annahme~\ref{ann:exc} lässt jedes gescheiterte \texttt{open} den Zustand unverändert (Lemma~\ref{lem:vertrag}(a)). Der Ablauf der Hooks ist dabei im Allgemeinen \emph{nicht} unverändert: Wurde eine Abhängigkeit neu geladen, bevor ein späterer Schritt scheiterte, so hat deren Init und anschließend deren Fini stattgefunden. (b) Ohne Annahme~\ref{ann:exc} gilt (a) nicht: Wirft der Init-Hook eines Handles eine andere Ausnahme als \texttt{SLError}, so bleiben die native Ressource abgebildet und die Referenzen seiner Abhängigkeiten erhöht.
\end{satz}

\begin{proof}
	(a) ist Lemma~\ref{lem:vertrag}(a); das Beispiel für die Hooks: Eine Datei $t$ mit Abhängigkeiten $b$ (gültig, neu) und $z$ (fehlt) lädt $b$, führt dessen Init aus, scheitert an $z$ und gibt $b$ frei, wobei dessen Fini läuft. (b) Der \texttt{except}-Zweig in \texttt{\_load} fängt nur \texttt{SLError}. Wirft \texttt{\_run\_init} etwa eine \texttt{OSError}, verlässt sie \texttt{\_load}, ohne dass \texttt{unmap} oder \texttt{\_release} der Abhängigkeiten ausgeführt wurden. Der Schlüssel wird zwar im \texttt{finally} vom Keller genommen, aber das Handle wurde nie eingetragen und seine Referenzen nie zurückgegeben; die Invariante aus Satz~\ref{satz:refcount} gilt dann nicht mehr. Dasselbe gilt für den Fini-Hook in \texttt{\_release}: Eine andere Ausnahme verlässt die Funktion, nachdem das Handle aus $\Lambda$ entfernt wurde und bevor \texttt{unmap} und die Rückgabe der Abhängigkeiten stattfinden.
\end{proof}

\section{Instanzen, Export und Nutzung}

\begin{satz}[Instanzidentität]
	\label{satz:instanz}
	Zwei erfolgreiche \texttt{open}-Aufrufe liefern genau dann dasselbe Handle-Objekt, wenn die zu ihnen gefundenen Dateien denselben Schlüssel (\texttt{realpath}) haben und das Handle dazwischen nicht freigegeben wurde. Insbesondere ist die Instanz durch die \emph{Datei}, nicht durch den logischen Namen bestimmt: Ein Name kann bei verschiedener Konfiguration oder verschiedenem Aufruferverzeichnis zu verschiedenen Dateien und damit zu verschiedenen gleichzeitigen Instanzen führen.
\end{satz}

\begin{proof}
	Der Schlüssel in $\Lambda$ ist \texttt{os.path.realpath(path)}; \texttt{\_open} liefert bei Treffer in $\Lambda$ das gespeicherte Objekt. Verschiedene Pfade auf dieselbe Datei (Symlinks) haben denselben \texttt{realpath}. Gegenbeispiel zur zweiten Aussage: Liegt \texttt{base.sl} in zwei Verzeichnissen $d_1,d_2$ und sind $d_1$, $d_2$ jeweils das Verzeichnis von zwei Bibliotheken $a$, $b$, die \texttt{base} deklarieren, und ist \texttt{base} in keinem Pfad vor dem Aufruferverzeichnis zu finden, dann lösen $a$ und $b$ \texttt{base} zu verschiedenen Dateien auf; beide Instanzen sind gleichzeitig geladen. Der Entwurf lässt offen, ob höchstens eine Instanz je Name erlaubt sein soll.
\end{proof}

\begin{satz}[Exportdurchsetzung]
	\label{satz:export}
	Sei $N(h)$ die Menge der Namen, für die \texttt{Loader.symbol}$(h,n)$ ohne Fehler zurückkehrt. Dann gilt $N(h)\subseteq X_h\cap S_h$, wobei $X_h$ die Exportliste in $M_h$ und $S_h$ die im nativen Payload auflösbaren Symbole sind. Gilt $X_h\subseteq S_h$, so ist $N(h)=X_h$. Die Durchsetzung ist eine Eigenschaft der API, keine Sicherheitsgrenze: Der Zugriff auf \texttt{handle.native.cdll} umgeht sie.
\end{satz}

\begin{proof}
	\texttt{symbol} wirft \texttt{SYMBOL\_NOT\_FOUND}, wenn $n\notin X_h$, und reicht sonst an \texttt{adapter.lookup} weiter, das für $n\notin S_h$ ebenfalls \texttt{SYMBOL\_NOT\_FOUND} wirft. Die Gleichheit folgt aus beiden Fällen. Das Feld \texttt{native} des öffentlichen Handles enthält das \texttt{ctypes}-Objekt, dessen Attributzugriff alle nativen Symbole auflöst. Die Hooks werden unabhängig von $X_h$ über \texttt{adapter.lookup} aufgerufen; sie müssen nicht exportiert sein.
\end{proof}

\begin{satz}[Nutzungssicherheit bei korrekter Klammerung]
	\label{satz:nutzung}
	Jeder Aufruf von \texttt{symbol} und jede Verwendung der zurückgegebenen Funktion durch einen Aufrufer liege zeitlich innerhalb des Intervalls zwischen dem \texttt{open}, das ihm das Handle lieferte, und dem zugehörigen \texttt{close} desselben Aufrufers, und kein Aufrufer rufe \texttt{close} zweimal für ein \texttt{open} auf. Dann ist das Handle während der Verwendung nicht freigegeben und nicht abgebildet.
\end{satz}

\begin{proof}
	Zu jedem Zeitpunkt der Verwendung ist $\mathrm{ext}(h)\ge1$, denn das \texttt{open} des Aufrufers ist noch nicht durch sein \texttt{close} beantwortet, und es wird nur einmal beantwortet. Nach Satz~\ref{satz:refcount} ist $\mathrm{rc}(h)\ge\mathrm{ext}(h)\ge1$, also $h\in\Lambda$ und $\texttt{closed}=\mathrm{falsch}$; das native Handle wird erst bei $\mathrm{rc}=0$ freigegeben. \texttt{open} und \texttt{close} sind durch die Sperre serialisiert, so dass die Zähler konsistent bleiben.
\end{proof}

Die Voraussetzung des Satzes ist wesentlich. Alle Aufrufer, die denselben Schlüssel öffnen, teilen sich \emph{ein} Handle-Objekt; ein doppeltes \texttt{close} eines Aufrufers nimmt deshalb die Referenz eines anderen zurück, und ohne die Klammerung kann \texttt{symbol} zwischen der Prüfung auf \texttt{closed} und dem Auflösen mit einem \texttt{close} eines anderen Threads zusammentreffen, weil \texttt{symbol} und \texttt{info} die Sperre nicht nehmen (Befund~\ref{bef:handle}). Die Klasse \texttt{binding.Library} verhindert das doppelte \texttt{close} durch Setzen ihres Feldes auf \texttt{None}.

\section{Nebenläufigkeit}

\begin{satz}[Deadlock durch die globale Sperre]
	\label{satz:deadlock}
	Es gibt ein Programm mit zwei Threads und einem gemeinsamen \texttt{Loader}, das nicht terminiert, obwohl keine der Bibliotheken zyklisch abhängt und keine Operation einen Fehler meldet.
\end{satz}

\begin{proof}
	\texttt{Loader.open} umschließt \texttt{\_open}, also auch den Init-Hook, mit \texttt{with self.\_lock} (Zeile~47). Sei $A$ eine Bibliothek, deren Init-Hook wartet, bis ein zweiter Thread ein Ereignis $\varepsilon$ gesetzt hat. Thread~1 ruft \texttt{open}($A$) auf; er hält die Sperre und blockiert im Hook. Thread~2 ruft zuerst \texttt{open}($B$) für eine beliebige Bibliothek $B$ auf und will danach $\varepsilon$ setzen. \texttt{open}($B$) verlangt die Sperre, die Thread~1 hält, und blockiert. Der Wartegraph hat die Kanten Thread~1~$\to$~Thread~2 (wartet auf $\varepsilon$) und Thread~2~$\to$~Thread~1 (wartet auf die Sperre): einen Kreis. Keiner der beiden kann fortfahren.
\end{proof}

Das verletzt die Zusage des Entwurfs (Abschnitt~5.7), Initialisierung laufe nie unter einer Sperre, die Nutzercode betreten kann; die README nennt die Sperre ausdrücklich (\emph{reentrant}, hält während des Initialisierens). Die Wiedereintrittsfähigkeit löst nur den Fall \emph{desselben} Threads: Ein \texttt{open} derselben Bibliothek im Init-Hook trifft auf den Schlüssel im Keller und endet in \texttt{CYCLE} (Satz~\ref{satz:erfolg}), nicht in einem Selbst-Deadlock. Abschnitt~\ref{sec:schluessel-sperren} des Ausblicks skizziert einen Entwurf mit Sperren je Schlüssel.

\chapter{Plattformadapter und Schichtung}
\label{chap:adapter}

\section{Der Adapter als Orakel}

Der Adapter ist eine abstrakte Klasse mit vier Operationen. Listing~\ref{lst:adapter} zeigt die gemeinsame Basisimplementierung, die den Payload in ein frisches temporäres Verzeichnis schreibt und mit \texttt{ctypes} öffnet.

\begin{lstlisting}[language=Python, caption={Abbilden und Freigeben (\texttt{adapters.py}, Zeilen 39--63 und 66--96, gekürzt).}, label={lst:adapter}]
def map_payload(self, payload: bytes, name: str) -> NativeLib:
    tmpdir = tempfile.mkdtemp(prefix="sl-")
    path = os.path.join(tmpdir, name + self.native_suffix)
    try:
        with open(path, "wb") as handle:
            handle.write(payload)
        return NativeLib(self._open(path), path, tmpdir)
    except OSError:
        shutil.rmtree(tmpdir, ignore_errors=True)
        raise

def unmap(self, native: NativeLib) -> None:
    self._release(native.cdll)
    shutil.rmtree(native.tmpdir, ignore_errors=True)

class LinuxAdapter(_PosixAdapter):    # system_dirs: /usr/local/lib/sl, /usr/lib/sl
class DarwinAdapter(_PosixAdapter):   # system_dirs: /usr/local/lib/sl, /opt/homebrew/lib/sl
class WindowsAdapter(PlatformAdapter):  # system_dirs: ProgramData\sl, ProgramFiles\sl
    def _release(self, cdll) -> None:
        _ctypes.FreeLibrary(cdll._handle)     # POSIX: _ctypes.dlclose(cdll._handle)
\end{lstlisting}

\begin{lemma}[Kollisionsfreie Abbildung]
	\label{lem:tmp}
	Zwei Aufrufe von \texttt{map\_payload} legen ihre Dateien in verschiedenen Verzeichnissen ab, auch bei gleichem Namen $\nu$. Die Datei wird freigegeben, bevor ihr Verzeichnis entfernt wird.
\end{lemma}

\begin{proof}
	\texttt{tempfile.mkdtemp} erzeugt ein bisher nicht vorhandenes Verzeichnis (nach der Dokumentation atomar, mit eindeutigem Namen). \texttt{unmap} ruft zuerst \texttt{\_release} auf (\texttt{dlclose} bzw. \texttt{FreeLibrary}) und danach \texttt{rmtree}; dies ist unter Windows wesentlich, weil geladene DLLs nicht gelöscht werden können.
\end{proof}

\begin{bemerkung}[Grenzen der Freigabe]
	\texttt{dlclose} gibt nach POSIX nur \emph{einen Verweis} auf das Objekt frei; ob das Abbild tatsächlich entladen wird, ist nicht zugesichert (etwa bei Objekten, die sich als nicht entladbar markiert haben). Die Referenzzählung des Laders (Satz~\ref{satz:refcount}) zählt Verweise des Laders, nicht Abbilder im Speicher. Die Abbildung über eine temporäre Datei setzt außerdem voraus, dass das Verzeichnis für temporäre Dateien Ausführung erlaubt; auf Systemen, die es ohne Ausführungsrecht einhängen, kann das Abbilden scheitern \offen{nicht geprüft}.
\end{bemerkung}

\section{Der Importgraph und die Schichten}

Die 14 Module von \texttt{sl} bilden einen Importgraphen. Er ist azyklisch und ordnet die Module in Ebenen (Abbildung~\ref{abb:importe}).

\begin{figure}[!ht]
	\centering
	\resizebox{0.97\textwidth}{!}{%
	\begin{tikzpicture}[font=\small, >={Stealth[length=2mm]},
		m/.style={draw=mainblue, fill=lightgray, rounded corners=2pt, minimum width=2.1cm, minimum height=0.7cm, inner sep=2pt}]
		\node[m] (er) at (0,0) {errors};
		\node[m] (pi) at (-1.5,1.5) {platform\_info};
		\node[m] (se) at (2.5,1.5) {search};
		\node[m] (me) at (-1.5,3.0) {metadata};
		\node[m] (ad) at (1.5,3.0) {adapters};
		\node[m] (co) at (-1.5,4.5) {container};
		\node[m] (lo) at (1.5,6.0) {loader};
		\node[m] (bu) at (-3.8,6.0) {build};
		\node[m] (ca) at (1.5,7.5) {capi};
		\node[m] (cg) at (-3.8,7.5) {codegen};
		\node[m] (cl) at (-1.2,9.0) {cli};
		\node[m] (bi) at (1.5,9.0) {binding};
		\node[m] (in) at (1.5,10.5) {\_\_init\_\_};
		\node[m] (mn) at (-1.2,10.5) {\_\_main\_\_};
		\draw[->] (pi)--(er); \draw[->] (se)--(er); \draw[->] (me)--(pi); \draw[->] (ad)--(pi); \draw[->] (co)--(me);
		\draw[->] (lo)--(co); \draw[->] (lo)--(ad); \draw[->] (lo)--(se); \draw[->] (bu)--(co); \draw[->] (ca)--(lo);
		\draw[->] (cg)--(bu); \draw[->] (cl)--(cg); \draw[->] (bi)--(ca); \draw[->] (in)--(bi); \draw[->] (mn)--(cl);
		\node[font=\scriptsize, anchor=west, text=darkgray] at (3.4,0) {Ebene 0};
		\node[font=\scriptsize, anchor=west, text=darkgray] at (3.9,1.5) {1};
		\node[font=\scriptsize, anchor=west, text=darkgray] at (3.4,3.0) {2};
		\node[font=\scriptsize, anchor=west, text=darkgray] at (3.4,4.5) {3};
		\node[font=\scriptsize, anchor=west, text=darkgray] at (3.4,6.0) {4};
		\node[font=\scriptsize, anchor=west, text=darkgray] at (3.4,7.5) {5};
		\node[font=\scriptsize, anchor=west, text=darkgray] at (3.4,9.0) {6};
		\node[font=\scriptsize, anchor=west, text=darkgray] at (3.4,10.5) {7};
	\end{tikzpicture}}
	\caption{Importgraph von \texttt{sl} (nur direkte Importe innerhalb des Pakets, jeweils zum höchsten Ziel gezeichnet; \texttt{loader} importiert zusätzlich \texttt{errors}, \texttt{metadata} und \texttt{platform\_info}, \texttt{build} importiert \texttt{errors}, \texttt{metadata} und \texttt{platform\_info}, \texttt{cli} importiert \texttt{build}, \texttt{codegen}, \texttt{container}, \texttt{errors}, \texttt{metadata}; \texttt{\_\_init\_\_} importiert \texttt{adapters}, \texttt{binding}, \texttt{capi}, \texttt{container}, \texttt{errors}, \texttt{loader}, \texttt{metadata}).}
	\label{abb:importe}
\end{figure}

\begin{satz}[Schichtenordnung und Laufzeitabschluss]
	\label{satz:schichten}
	Der Importgraph von \texttt{sl} ist azyklisch. Die Abbildung $\lambda$ mit $\lambda(\texttt{errors})=0$; $\lambda(\texttt{platform\_info})=\lambda(\texttt{search})=1$; $\lambda(\texttt{metadata})=\lambda(\texttt{adapters})=2$; $\lambda(\texttt{container})=3$; $\lambda(\texttt{loader})=\lambda(\texttt{build})=4$; $\lambda(\texttt{capi})=\lambda(\texttt{codegen})=5$; $\lambda(\texttt{cli})=\lambda(\texttt{binding})=6$; $\lambda(\texttt{\_\_init\_\_})=\lambda(\texttt{\_\_main\_\_})=7$ erfüllt $\lambda(u)>\lambda(v)$ für jeden Import $u\to v$. Der Abschluss von \texttt{loader} unter Importen ist $\{\texttt{loader},\texttt{adapters},\texttt{container},\texttt{errors},\texttt{metadata},\texttt{platform\_info},\texttt{search}\}$; er enthält weder \texttt{build}, \texttt{codegen}, \texttt{cli}, \texttt{capi} noch \texttt{binding}.
\end{satz}

\begin{proof}
	Die Importe sind den \texttt{import}-Zeilen der Module zu entnehmen: \texttt{platform\_info}$\to$\texttt{errors}; \texttt{search}$\to$\texttt{errors}; \texttt{metadata}$\to$\texttt{errors}, \texttt{platform\_info}; \texttt{adapters}$\to$\texttt{errors}, \texttt{platform\_info}; \texttt{container}$\to$\texttt{errors}, \texttt{metadata}; \texttt{loader}$\to$\texttt{adapters}, \texttt{container}, \texttt{errors}, \texttt{metadata}, \texttt{platform\_info}, \texttt{search}; \texttt{build}$\to$\texttt{container}, \texttt{errors}, \texttt{metadata}, \texttt{platform\_info}; \texttt{capi}$\to$\texttt{errors}, \texttt{loader}; \texttt{codegen}$\to$\texttt{build}, \texttt{metadata}; \texttt{binding}$\to$\texttt{capi}; \texttt{cli}$\to$\texttt{build}, \texttt{codegen}, \texttt{container}, \texttt{errors}, \texttt{metadata}; \texttt{\_\_init\_\_}$\to$\texttt{adapters}, \texttt{binding}, \texttt{capi}, \texttt{container}, \texttt{errors}, \texttt{loader}, \texttt{metadata}; \texttt{\_\_main\_\_}$\to$\texttt{cli}. Für jeden dieser Importe $u\to v$ ist $\lambda(u)=1+\max\lambda(v)$ über die Ziele $v$ von $u$ (nachzurechnen: \texttt{loader}: $1+\max(2,3,0,2,1,1)=4$; \texttt{build}: $1+\max(3,0,2,1)=4$; \texttt{capi}: $1+4=5$; \texttt{codegen}: $1+\max(4,2)=5$; \texttt{binding}: $1+5=6$; \texttt{cli}: $1+\max(4,5,3,0,2)=6$; \texttt{\_\_init\_\_}: $1+\max(2,6,5,3,0,4,2)=7$; \texttt{\_\_main\_\_}: $1+6=7$). Eine streng fallende Ebenenfunktion schließt Kreise aus. Der Abschluss von \texttt{loader} ergibt sich durch Vereinigung der Ziele: \texttt{adapters}, \texttt{container}, \texttt{errors}, \texttt{metadata}, \texttt{platform\_info}, \texttt{search} und deren Ziele, die in dieser Menge liegen.
\end{proof}

Praktisch heißt das: Wer nur laden will, braucht weder Compiler-Aufrufe noch Codegenerierung noch Kommandozeile; der Laufzeitpfad ist ein klar abgegrenzter Teil des Pakets. Das Paket hat keine Laufzeitabhängigkeiten außerhalb der Standardbibliothek (\texttt{dependencies = []}).

\section{Wie weit trägt die Schichtung?}

\begin{satz}[Adapter-Parametrizität mit Ausnahmen]
	\label{satz:parametrisch}
	Die Beobachtung einer Operationsfolge des Laders hängt vom Betriebssystem nur ab über
	\begin{enumerate}[label=(\arabic*)]
		\item die vier Operationen des Adapters (\texttt{system\_dirs}, \texttt{map\_payload}, \texttt{lookup}, \texttt{unmap}),
		\item die Pfadfunktionen der Standardbibliothek \texttt{os.path.isabs}, \texttt{join}, \texttt{realpath}, \texttt{dirname}, \texttt{abspath}, \texttt{isfile}, \texttt{basename},
		\item die Verzweigung \texttt{os.name == "nt"} in \texttt{search.\_world\_writable},
		\item die Vorgabewerte \texttt{current\_platform()} und \texttt{current\_arch()} bei der Konstruktion eines \texttt{Loader} ohne ausdrückliche Angabe.
	\end{enumerate}
	Mit Fake-Adapter, expliziter Plattform und Architektur und einer POSIX-Pfadumgebung ist die Beobachtung eine Funktion von $(E,R,A,h,\mathcal{F},\mathcal{O})$.
\end{satz}

\begin{proof}
	Wir durchsuchen die Module des Laufzeitpfads (Satz~\ref{satz:schichten}) nach Zugriffen auf Plattformzustand. \texttt{loader.py}: \texttt{os.environ} nur als Vorgabe für \texttt{env}; \texttt{os.path.realpath} und \texttt{dirname}; die Vorgaben aus \texttt{platform\_info}; Orakelaufrufe nur über \texttt{self.adapter} an den Stellen \texttt{system\_dirs} (über \texttt{locate}), \texttt{map\_payload}, \texttt{lookup} (Symbole, Init, Fini) und \texttt{unmap}. \texttt{search.py}: \texttt{os.path.isabs}, \texttt{join}, \texttt{isfile}, \texttt{abspath}, \texttt{dirname}, \texttt{basename}, \texttt{os.stat}, \texttt{os.name}. \texttt{container.py} und \texttt{metadata.py} greifen auf keine Plattformfunktion zu; \texttt{platform\_info.py} liest \texttt{sys.platform} und \texttt{platform.machine()} nur in den Vorgabewerten. Alle weiteren Aufrufe sind reine Funktionen ihrer Argumente. Die Zusatzaussage folgt, weil unter diesen Annahmen (1), (3) und (4) entfallen und (2) unter POSIX-Pfaden eine feste Funktion ist.
\end{proof}

Die Tabelle~\ref{tab:leaks} stellt die Ausnahmen den Orten im Quelltext gegenüber. Der Entwurf sagt, \emph{alles} Plattformspezifische liege im Adapter; die Referenzimplementierung hält dies für die Funktionen des Systemladers ein, nicht für Pfadsyntax und Rechteprüfung. Das ist kein Mangel an Sorgfalt, sondern folgt aus der Wahl, Pfade als Zeichenketten zu behandeln; der Ausblick (Abschnitt~\ref{sec:adapter-ausblick}) beschreibt eine saubere Trennung.

\begin{table}[!ht]
	\centering
	\small
	\begin{tabularx}{\textwidth}{@{}L{3.1cm}L{4.3cm}Y@{}}
		\toprule
		Ort & Plattformzugriff & Folge \\
		\midrule
		\texttt{search.py}:84, 94, 137, 138, 143 & \texttt{os.path.isabs}, \texttt{join}, \texttt{isfile}, \texttt{dirname}, \texttt{abspath} & Pfadsyntax (\texttt{/} oder \texttt{\textbackslash}) und Absolutheit hängen vom System ab; derselbe \texttt{SL\_PATH}-Eintrag \texttt{C:\textbackslash x} ist auf Linux relativ und wird verworfen \\[0.3em]
		\texttt{search.py}:120--127 & \texttt{os.name}, \texttt{os.stat} & Rechteprüfung nur auf POSIX; unter Windows Testnamen-Sonderfall (B3) \\[0.3em]
		\texttt{loader.py}:87, 146 & \texttt{os.path.realpath}, \texttt{dirname} & Instanzidentität (Satz~\ref{satz:instanz}) hängt von der Symlink-Auflösung des Systems ab \\[0.3em]
		\texttt{loader.py}:33--35 & \texttt{current\_platform}, \texttt{current\_arch} & Vorgaben des Laders \\[0.3em]
		\texttt{adapters.py}:43, 44, 90 & \texttt{tempfile.mkdtemp}, \texttt{os.path.join} & Ort der Materialisierung hängt von \texttt{TMPDIR}/\texttt{TEMP} ab \\
		\bottomrule
	\end{tabularx}
	\caption{Plattformzugriffe außerhalb der vier Adapteroperationen (Zeilennummern des hochgeladenen Standes).}
	\label{tab:leaks}
\end{table}

\begin{bemerkung}[Was die Fake-Adapter-Tests beweisen]
	Die Tests des Lader-Kerns laufen mit einem \texttt{FakeAdapter}, der keine native Bibliothek lädt (\texttt{tests/conftest.py}). Nach Satz~\ref{satz:parametrisch} ist die Beobachtung des Lader-Kerns unter diesen Bedingungen unabhängig vom Betriebssystem. Ein Lauf dieser Tests auf einem System ist damit ein Lauf der Lader-Logik für \emph{alle} Systeme, soweit nicht die Ausnahmen (2)--(4) berührt werden. Die Tests der echten Adapter für Windows und macOS laufen über Nachbildungen (README), nicht auf diesen Systemen.
\end{bemerkung}

